\chapter{Reproducibility}
\label{sec:repro}

\section{Reproducing the Addition Arithmetic Circuit}

\citet{lindsey2025biology} showed that Claude~3.5 Haiku, when presented
with prompts of the form \texttt{calc:~$a$+$b$=}, does not implement a
symbolic carry algorithm but instead activates a family of sparse
\textit{lookup table} features whose firing concentrates on the cells of
the operand space satisfying $a \bmod 10 = d_1$ and $b \bmod 10 = d_2$
for a specific digit pair $(d_1, d_2)$.  These features encode the ones
digit of $d_1 + d_2$ and route this information directly to the output.
The present reproduction asks whether the same mechanistic signature
appears in Qwen3-4B, trained independently on a different corpus and with
a different tokeniser, and whether the attribution methodology transfers
faithfully to a model for which it was not originally calibrated.

\subsection{Transcoders and the Linearised Computation Graph}

Let $\mathcal{M}$ be a transformer with $L$ layers, residual stream
dimension $d$, and MLP sublayers $f_0, \ldots, f_{L-1}$.  Mechanistic
analysis replaces each $f_\ell$ with a pretrained \textit{transcoder}
$\mathcal{T}_\ell$, a sparse encoder--decoder whose weights are fit
offline to reconstruct the MLP output while decomposing it into $K$
interpretable directions.  For a residual-stream vector
$\mathbf{x} \in \mathbb{R}^d$ entering layer $\ell$, the transcoder
computes pre-activations

\begin{equation}
  z_{\ell,k}(\mathbf{x})
    = \bigl(\mathbf{W}_\mathrm{enc}^{(\ell)}\,\mathbf{x}
           + \mathbf{b}_\mathrm{enc}^{(\ell)}\bigr)_k,
  \label{eq:preact}
\end{equation}

so that feature $k$ carries activation
$a_{\ell,k}(\mathbf{x}) = \mathrm{ReLU}(z_{\ell,k}(\mathbf{x}))$,
and the reconstructed MLP output is

\begin{equation}
  \hat{f}_\ell(\mathbf{x})
    = \mathbf{W}_\mathrm{dec}^{(\ell)}\,\mathbf{a}_\ell(\mathbf{x})
    + \mathbf{b}_\mathrm{dec}^{(\ell)},
  \label{eq:recon}
\end{equation}

where $\mathbf{W}_\mathrm{dec}^{(\ell)} \in \mathbb{R}^{d \times K}$ and
$\mathbf{a}_\ell(\mathbf{x})$ is the vector of all feature activations at
layer $\ell$.  Throughout this section a double subscript $(\ell, k)$
indexes features, with $\ell$ specifying the layer and $k$ the index within
that layer's dictionary, in place of the flat index $f$ used in
Chapter~\ref{sec:introduction}.

Attribution is computed on the \textit{linearised} model, in which attention
outputs are detached and treated as constants, so the residual stream
propagates solely through the transcoder reconstruction path.  This
linearisation is the necessary condition for the attribution scores to
be additive and interpretable as influence flows; it is also what
constrains their faithfulness, a point examined in Section~\ref{sec:faithfulness}.

Let $T$ denote the position of the final input token, and let
$y_j = \mathbf{v}_j^\top \mathbf{h}_T$ be the logit for output token
$j$, where $\mathbf{h}_T$ is the pre-unembedding hidden state and
$\mathbf{v}_j$ is the (demeaned) unembedding direction for token $j$.
The \textit{attribution score} of feature $k$ at position $t$ in
layer $\ell$ is defined as

\begin{equation}
  \alpha_{\ell,k,t}
    \;=\; a_{\ell,k,t}
          \cdot \frac{\partial y_j}{\partial a_{\ell,k,t}},
  \label{eq:attr}
\end{equation}

the product of the feature activation $a_{\ell,k,t}$ with the gradient of
the target logit with respect to that activation, consistent with
Equation~\eqref{eq:attribution} of Chapter~\ref{sec:introduction}.
Collecting all such scores over a given prompt yields a directed weighted
graph, the \textit{attribution graph}, in which nodes are (layer, position,
feature) triples and edges carry the influence $\alpha_{\ell,k,t}$ from
earlier features to later ones.  The graph is then pruned by retaining only the nodes and edges
that account for a fixed fraction of the total downstream influence on
$y_j$, controlled by a node threshold $\tau_\mathrm{node}$ and an edge
threshold $\tau_\mathrm{edge}$.

\subsection{Operand Plots and Lookup Feature Identification}

For a prompt $\texttt{calc:~}a\texttt{+}b\texttt{=~}$, let $T$ be the
index of the \texttt{=} token.  The \textit{operand matrix} for feature
$(\ell, k)$ is

\begin{equation}
  M_{\ell,k}[a, b]
    \;=\; a_{\ell,k,T}
          \!\left(\texttt{calc:~}a\texttt{+}b\texttt{=~}\right),
    \qquad a, b \in \{0, \ldots, 99\},
  \label{eq:operand-matrix}
\end{equation}

obtained by sweeping over the full $100 \times 100$ integer-pair grid.
The geometry of $M_{\ell,k}$ encodes the feature's sensitivity in a
directly interpretable form.  A feature that tracks the raw sum $a + b$
produces diagonal stripes of constant activation; one that tracks a
single operand produces horizontal or vertical bands; modular sensitivity
to the ones digit of either operand produces a period-10 repeating
structure.  A feature that fires selectively for a specific pair of ones
digits $(d_1, d_2)$ produces a regular grid of isolated point clusters,
concentrated on the $200$ cells where $a \bmod 10 \in \{d_1, d_2\}$ and
$b \bmod 10 \in \{d_1, d_2\}$.  These \textit{lookup table features}
constitute $2\%$ of the operand space; a uniform feature would achieve
the same coverage by chance.

To identify lookup candidates automatically, each feature is assigned the
concentration score

\begin{equation}
  s_{\ell,k}(d_1, d_2)
    \;=\;
    \frac{%
      \displaystyle\sum_{a,b}
        M_{\ell,k}[a,b]\;
        \mathbb{1}\!\left[
          a \bmod 10 \in \{d_1, d_2\},\;
          b \bmod 10 \in \{d_1, d_2\}
        \right]
    }{%
      \displaystyle\sum_{a,b} M_{\ell,k}[a,b]
    },
  \label{eq:lookup-score}
\end{equation}

which measures the fraction of total activation mass that lies on the
target grid.  Features are filtered by a minimum peak-activation
threshold to exclude dormant features, then ranked by $s_{\ell,k}$.
The top-ranked feature for the pair $(d_1, d_2) = (6, 9)$ -- the ones
digits of the focus operands $36$ and $59$ -- is taken as the reference
lookup feature and subjected to causal validation.  Two interventions are
applied.  In the \textit{inhibition} experiment, the feature's activation
is scaled by a factor $\beta \in \{1.0, 0.5, 0.0, -1.0, -5.0, -10.0\}$
before the decoder output is computed; if the feature is causally
necessary, driving $\beta$ below zero should monotonically weaken the
logit for the token representing the correct ones digit of $6 + 9 = 15$.
In the \textit{substitution} experiment, the MLP inputs at all layers
preceding the lookup feature's layer are replaced by those from the
prompt \texttt{calc:~39+59=~}, which has ones digits $(9, 9)$ rather than
$(6, 9)$; if the identified feature mediates the computation, the output
should shift from predicting ones digit $5$ to ones digit $8$.

\subsection{Experimental Scope}
\label{subsec:addition-dataset}

The experiment is structured around four prompt collections,
summarised in Table~\ref{tab:addition-dataset}.

\begin{table}[ht]
  \centering
  \small
  \caption{%
    Prompt collections used in the addition circuit reproduction.
    The operand grid drives the activation sweep; the focus prompt
    anchors the attribution graph.
  }
  \label{tab:addition-dataset}
  \begin{tabular}{@{}llrl@{}}
    \toprule
    Split & Template & $N$ & Purpose \\
    \midrule
    Operand grid
      & \texttt{calc:~$a$+$b$=~}
      & $10{,}000$
      & Feature activation sweep; $a,b \in \{0,\ldots,99\}$ \\
    Focus prompt
      & \texttt{calc:~36+59=~}
      & $1$
      & Primary attribution graph; answer \texttt{95} \\
    \bottomrule
  \end{tabular}
\end{table}

The $10{,}000$-prompt grid exhausts the integer-pair space
$\{0,\ldots,99\}^2$, covering all $100$ distinct ones-digit
combinations and providing sufficient support to identify modular
activation patterns visually.  The focus prompt \texttt{calc:~36+59=}
follows Anthropic's exact choice, enabling direct comparison of
attribution graphs across models.

\subsection{Behavioural Verification}
\label{subsec:accuracy-sweep}

Before committing mechanistic resources to a single focus prompt, it is
necessary to establish that the model solves the task reliably across the
operand space.  A circuit identified on a prompt the model cannot answer
carries no mechanistic authority.  The accuracy sweep therefore records,
for every pair $a, b \in \{0,\ldots,99\}$, the \textit{teacher-forced
confidence} at each output position.  Let
$\mathbf{y}^*(a,b) = (y^*_0, y^*_1, \ldots)$ denote the ground-truth
answer tokens for $a + b$.  The teacher-forced confidence at position $k$
is

\begin{equation}
  P_k(a, b)
  = P\!\left(
      y_k = y^*_k \;\middle|\;
      y_0 = y^*_0,\, \ldots,\, y_{k-1} = y^*_{k-1},\;
      \mathbf{x}(a,b)
    \right),
  \label{eq:teacher-forced}
\end{equation}

where $\mathbf{x}(a,b)$ is the token sequence for the prompt
\texttt{calc:~$a$+$b$=} and preceding positions are forced to their
correct tokens.  Equation~\eqref{eq:teacher-forced} isolates the residual
difficulty at each generation step by removing uncertainty due to
error accumulation in prior positions.

Figure~\ref{fig:positional-cascade} plots $P_k(a,b)$ for
$k \in \{0, 1, 2\}$.  At $k = 0$ the model must predict the leading
output token from the prompt alone; the grid mean is $\mu_0 = 0.233$,
reflecting genuine uncertainty over a substantial fraction of the operand
space.  At $k = 1$, once the first token is teacher-forced to the correct
value, mean confidence rises to $\mu_1 = 0.943$.  By $k = 2$ it reaches
$\mu_2 = 0.988$, with the only uncertain region being the upper-right
triangle where $a + b \geq 100$ and the answer carries into a third digit.
The cascade reveals that the central computational burden falls entirely on
predicting $y^*_0$: subsequent tokens resolve almost automatically given
the correct leading digit, a pattern that is consistent with a mechanism
that maps directly from the operand ones-digits to the output rather than
propagating carries sequentially.

\begin{figure}[ht]
  \centering
  \includegraphics[width=\textwidth]{../runs/addition/accuracy_sweep/plots/positional_cascade}
  \caption{%
    Teacher-forced confidence $P_k(a,b)$ (Eq.~\ref{eq:teacher-forced})
    across the $100 \times 100$ operand grid for positions $k = 0, 1, 2$.
    Each cell encodes the probability that Qwen3-4B assigns to the correct
    $k$-th output token when all preceding tokens are forced to their
    ground-truth values.  Grid means are $\mu_0 = 0.233$, $\mu_1 = 0.943$,
    and $\mu_2 = 0.988$.  The leftmost panel reveals a period-10 grid structure
    aligned with the decimal ones-digit boundaries: columns and rows separated
    by ten units carry systematically different baseline confidence levels,
    implying that $P_0$ depends primarily on $(a \bmod 10,\, b \bmod 10)$
    rather than on operand magnitudes.  Empty cells in the middle panel
    correspond to pairs with $a + b < 10$, whose single-digit sum has no
    second output token; empty cells in the rightmost panel correspond to
    pairs with $a + b < 100$, whose two-digit sum has no third digit.
  }
  \label{fig:positional-cascade}
\end{figure}

Rather than varying smoothly with operand magnitude, the $k = 0$ panel of
Figure~\ref{fig:positional-cascade} shows that confidence is organised into
a period-10 grid aligned with the decimal ones-digit boundaries.  Within
each $10 \times 10$ block the within-block pattern repeats with the same
geometry, and this periodicity implies that $P_0$ depends primarily on
$(a \bmod 10,\, b \bmod 10)$ rather than on the absolute values of $a$ and
$b$.  This is the behavioural signature one would expect if the computation
were implemented by features selective to specific ones-digit pairs.

\begin{figure}[ht]
  \centering
  \includegraphics[width=0.62\textwidth]{../runs/addition/accuracy_sweep/plots/entropy_map}
  \caption{%
    Prediction entropy $H(a,b)$ (Eq.~\ref{eq:entropy}) at position $k = 0$.
    Low-entropy cells (blue) form a period-10 grid aligned with the
    high-confidence cells visible in the $k = 0$ panel of
    Figure~\ref{fig:positional-cascade}, indicating that the model is both
    decisive and correct on precisely the pairs where lookup features activate
    strongly.  High-entropy cells (red) concentrate near carry boundaries and
    for large operand values, reflecting genuine distributional indecision
    rather than merely confident errors.
  }
  \label{fig:entropy-map}
\end{figure}

\paragraph{Prediction entropy.}
Confidence $P_k(a,b)$ measures whether the top-ranked prediction is
correct, but not how sharply the model has committed to any prediction.
A cell with low confidence might correspond either to genuine
distributional uncertainty or to a confident wrong prediction.  The
Shannon entropy of the truncated output distribution at position $0$,

\begin{equation}
  H(a, b)
  = -\sum_{j=1}^{K} p_{(j)}(a, b)\;\log_2 p_{(j)}(a, b),
  \label{eq:entropy}
\end{equation}

where $p_{(1)}(a,b) \geq \cdots \geq p_{(K)}(a,b)$ are the top-$K$
predicted probabilities, measures distributional spread independently of
prediction correctness.  Figure~\ref{fig:entropy-map} plots $H(a,b)$
across the grid.  Low-entropy cells (where the model commits decisively to
a single token) form a period-10 grid aligned precisely with the
high-confidence regions of the $k = 0$ panel of
Figure~\ref{fig:positional-cascade}.
High-entropy cells (where the model distributes probability across several
candidates) are concentrated at and near carry boundaries and for large
operand values.  The agreement between the two maps establishes that the
model's decisiveness and its correctness are driven by the same underlying
structure, rather than by one compensating for the other.

\paragraph{Carry complexity.}
Each pair $(a,b)$ can be classified by the number of digit-column carries
arising in the sum.  Defining

\begin{equation}
  \kappa(a, b)
  = \begin{cases}
      0 & \text{no carry arises,} \\
      1 & \text{exactly one carry arises,} \\
      2 & \text{two or more carries arise,}
    \end{cases}
  \label{eq:carry-class}
\end{equation}

Figure~\ref{fig:carry-structure} overlays the partition induced by
$\kappa$ on the confidence grid.  If Qwen3-4B were implementing a symbolic
carry algorithm, accuracy should deteriorate with carry complexity:
pairs with $\kappa = 2$ should be hardest and pairs with $\kappa = 0$
easiest.  The right panel of Figure~\ref{fig:carry-structure} contradicts
this prediction.  Regions of identical carry class exhibit substantial
internal variation in $P_0$, and the drops in confidence cut across carry
boundaries rather than aligning with them.  The period-10 grid structure
of $P_0$ persists within each carry class, unmodulated by carry count.
This dissociation between carry complexity and model accuracy constitutes
behavioural evidence that the computation is not organised around carry
propagation.

\begin{figure}[ht]
  \centering
  \includegraphics[width=\textwidth]{../runs/addition/accuracy_sweep/plots/carry_structure}
  \caption{%
    Left: carry-class partition $\kappa(a,b)$
    (Eq.~\ref{eq:carry-class}) of the operand grid, coloured by carry count.
    Right: first-token confidence $P_0(a,b)$ with carry-boundary contours
    superimposed.  The period-10 structure of $P_0$ persists within each
    carry class, and drops in confidence do not align with the
    no-carry/single-carry/multi-carry boundaries, ruling out a symbolic
    carry algorithm as the underlying mechanism.
  }
  \label{fig:carry-structure}
\end{figure}

Taken together, the three diagnostics converge on the same conclusion.
The teacher-forced cascade localises the computation to the first output
token and reveals that the model's confidence is determined almost entirely
at that step.  The confidence and entropy maps show that this determination
is modular in the decimal ones digits of both operands, with a period-10
structure that would be inexplicable under a positional carry algorithm.
The carry-class overlay confirms that carry complexity does not predict
accuracy, removing the most natural alternative hypothesis.  These
behavioural observations are consistent with the lookup-table mechanism
reported by \citet{lindsey2025biology}, and they motivate the mechanistic
analysis of the following sections.
